% % =========================
% % report/results.tex
% % =========================

% % \subsection{نمایش کیفی هم‌پوشانی‌ها}
% % در \cref{fig:overlay_grid} نمونه‌هایی از خروجی‌های هم‌پوشانی ارائه شده است. هدف از این شکل، ارزیابی کیفی «درست نشستن خطوط/مرزهای الگو» روی ساختار متناظر در عکس است. اگر هم‌ترازی درست باشد، خطوط الگو باید با لبه‌ها و ساختارهای اصلی قطعه هم‌خوانی داشته باشند.

% % \ResultGraphic{figures/overlay_grid.png}{0.95\linewidth}{شبکه‌ای از نمونه‌های هم‌پوشانی الگو بر روی عکس.}
% % \label{fig:overlay_grid}

% \subsection{نتایج کمی (خلاصهٔ معیارها)}
% در جدول \cref{tab:summary_metrics} خلاصهٔ معیارهای کمی ذخیره‌شده در \texttt{summary.csv} گزارش می‌شود. اگر فایل موجود نباشد، گزارش به‌جای خطا، یک کادر «فایل پیدا نشد» نمایش می‌دهد (برای جلوگیری از شکست کامپایل).

% \SafeCSVTable{../results/metrics/summary.csv}{خلاصهٔ معیارهای سطح داده (از \texttt{summary.csv}).}{tab:summary_metrics}

% \subsection{توزیع خطای بازفرافکنی}
% هیستوگرام خطای بازفرافکنی (در صورت وجود زمین‌حقیقت) در \cref{fig:reproj_hist} نمایش داده می‌شود. وجود دنبالهٔ بلند در هیستوگرام معمولاً نشان‌دهندهٔ شکست‌های جدی (مثلاً هموگرافی اشتباه ولی سازگار) است.

% \ResultGraphic{figures/reproj_error_hist.png}{0.85\linewidth}{هیستوگرام خطای بازفرافکنی (اگر زمین‌حقیقت موجود است).}
% \label{fig:reproj_hist}

% \subsection{تحلیل زمان اجرا}
% نمودار تفکیک زمان مراحل مختلف در \cref{fig:runtime_breakdown} نشان می‌دهد کدام بخش گلوگاه است. در کاربرد واقعی، معمولاً استخراج ویژگی و تطبیق بیشترین سهم را دارند (بسته به تعداد ویژگی‌ها و اندازه تصویر).

% \ResultGraphic{figures/runtime_breakdown.png}{0.85\linewidth}{تفکیک زمان اجرا برای مراحل اصلی.}
% \label{fig:runtime_breakdown}

% \subsection{نتایج ابلیشن}
% در \cref{fig:ablation_plot} حساسیت سامانه به پارامترها نمایش داده می‌شود. هدف از ابلیشن پیدا کردن «ناحیهٔ عملکرد پایدار» است؛ یعنی تنظیماتی که با تغییرات کوچک پارامتر، افت شدید ایجاد نمی‌کند.

% \ResultGraphic{figures/ablation_plot.png}{0.85\linewidth}{نتایج ابلیشن: حساسیت عملکرد نسبت به پارامترها.}
% \label{fig:ablation_plot}

% \subsection{نمونه‌های شکست و تحلیل مکانیزم}
% شکل‌های \cref{fig:failure1,fig:failure2} نمونه‌هایی از شکست را نشان می‌دهند. شکست‌ها را می‌توان به چند کلاس رایج تقسیم کرد:
% \begin{itemize}
%   \item \textbf{کمبود ویژگی معتبر:} تعداد نقاط کلیدی یا تطبیق‌ها کم است.
%   \item \textbf{تطبیق‌های غلط ولی منسجم:} پس‌زمینه یا الگوهای تکراری باعث هموگرافی اشتباه می‌شود.
%   \item \textbf{نقض فرض صفحه‌ای:} سطح سه‌بعدی/خمیده باعث می‌شود هموگرافی یکپارچه کافی نباشد.
% \end{itemize}

% \ResultGraphic{figures/failure_case_1.png}{0.85\linewidth}{نمونه شکست ۱ (علت را با خروجی‌های دیباگ تحلیل کنید).}
% \label{fig:failure1}

% \ResultGraphic{figures/failure_case_2.png}{0.85\linewidth}{نمونه شکست ۲ (علت را با خروجی‌های دیباگ تحلیل کنید).}
% \label{fig:failure2}
